\chapter{Text shaping and combining}
\label{ch:shaping}

\begin{aside}
``Shape'' a string means turn code points into positioned
glyphs. Browsers do this with HarfBuzz; terminals mostly don't
do it at all. A monospace grid is simpler --- but only mostly.
\end{aside}

\section{Why shaping}

In proportional fonts, shaping decides which glyph to use for
each code point (small caps, swashes, kerning pairs).
Monospace fonts mostly skip this --- each code point maps to
one glyph at one width.

But some scripts need shaping even in monospace:

\begin{itemize}[leftmargin=*]
  \item \textbf{Arabic.} Letters have initial/medial/final
    forms; the same code point picks different glyphs based
    on context.
  \item \textbf{Devanagari.} Consonant clusters merge into
    conjuncts.
  \item \textbf{Combining marks.} Stacking diacritics.
\end{itemize}

Terminals delegate these to their font and renderer. \maya{}
sees ``code points''; the terminal sees ``glyphs''.

\section{What \maya{} does and doesn't}

\maya{} is responsible for:

\begin{itemize}[leftmargin=*]
  \item Grouping code points into grapheme clusters.
  \item Computing each cluster's column width.
  \item Storing the base code point in the canvas cell and
    combining marks in a side-table.
\end{itemize}

\maya{} is \emph{not} responsible for:

\begin{itemize}[leftmargin=*]
  \item Selecting glyph variants.
  \item Drawing combining marks at the right offset.
  \item Substituting ligatures (\code{->}, \code{!=}, etc.).
  \item Handling bidirectional text.
\end{itemize}

These belong to the terminal emulator's font renderer.

\section{Ligatures}

Programmer fonts like Fira Code substitute \code{->} with an
arrow glyph spanning two cells. The terminal sees two code
points and one cell each; the font renderer decides to draw
one arrow across two cells.

The diff cannot tell. If a ligature changes mid-frame, the
terminal redraws the relevant cells; both cells appear in the
diff.

Some users find ligatures confusing in TUIs (the cursor
position becomes ambiguous). \maya{} does not opt in or out
--- the font picks.

\section{Combining marks in the canvas}

A cell holds a base code point. Combining marks for that cell
live in a side-table keyed by cell index:

\begin{lstlisting}
struct CellExtras {
    std::vector<char32_t> marks;
};
std::unordered_map<size_t, CellExtras> extras_;
\end{lstlisting}

For most cells, no marks: no side-table entry. Cells with
marks pay a small allocation cost.

\section{Emitting combining marks}

When emitting a cell with combining marks, \maya{} emits the
base code point followed by each mark:

\begin{lstlisting}
emit_codepoint(base);
for (auto m : marks) emit_codepoint(m);
\end{lstlisting}

The terminal stacks the marks visually. Most monospace fonts
have decent combining-mark coverage; some scripts (Hebrew with
nikud) look broken on some fonts.

\section{ZWJ sequences}

Zero-width joiner sequences (the family emoji, the trans
flag, the rainbow flag) are emitted code point by code point;
the terminal recognises the sequence and draws one glyph.

\maya{} treats them as one grapheme cluster. The canvas
stores the first code point as base and the rest as
``combining'' (even though ZWJ is not technically a
combining mark; the model is the same).

\section{Variation selectors}

\code{U+FE0F} (emoji presentation) and \code{U+FE0E} (text
presentation) are stored as marks on the base. The terminal
honours them when drawing.

\section{Width and rendering interaction}

The width of a cluster determines column reservation; the
rendering decides what pixels appear in those columns. They
should agree. They don't always:

\begin{itemize}[leftmargin=*]
  \item A wide emoji is reserved 2 columns but a font might
    only draw 1 column wide.
  \item A combining mark is width 0 but the font might draw
    it overflowing into the next column.
\end{itemize}

\maya{} cannot detect these mismatches; the user sees a
visual artefact. The fix is a better font or a different
terminal.

\begin{warning}
If your app uses many emoji and combining marks, test on
multiple terminals and fonts. The variance is large.
\end{warning}

\section{Bidirectional text (out of scope)}

\maya{} renders code points left-to-right in memory order.
For RTL scripts, this is wrong but not catastrophic --- the
characters appear in reverse. The user can read them; the
caret semantics are confusing.

If you need BiDi, reverse RTL substrings yourself before
passing to \maya{}, or wait for a future version with
\code{libfribidi} integration.

\section{Recap}

\begin{recap}
\begin{itemize}[leftmargin=*]
  \item \maya{} clusters but does not shape.
  \item Combining marks stored in a side-table.
  \item ZWJ sequences are clusters; the terminal renders.
  \item Ligatures are font features; \maya{} is unaware.
  \item Variation selectors flip emoji/text presentation.
  \item BiDi is currently a user problem.
\end{itemize}
\end{recap}

\section*{Workshop}

\begin{enumerate}[leftmargin=*]
  \item Type Hebrew text and observe the order; reverse it
    in your code and confirm visual correctness.
  \item Compare a ZWJ family emoji rendered in your app to
    rendered in your terminal alone; should match.
  \item Use a font with ligatures; observe \code{->}
    becoming an arrow without \maya{} knowing.
\end{enumerate}
